% Standalone problem document, generated from the adjacent README.md.
% Compile with XeLaTeX. Mathematical content is maintained in Markdown.
\documentclass[11pt,a4paper]{article}
\usepackage[fontsize=10.5pt]{scrextend}
\usepackage[margin=22mm,headheight=15pt,headsep=9mm,footskip=11mm]{geometry}
\usepackage{fontspec}
\setmainfont{texgyrepagella}[Extension=.otf,UprightFont=*-regular,BoldFont=*-bold,ItalicFont=*-italic,BoldItalicFont=*-bolditalic]
\setsansfont{texgyreheros}[Extension=.otf,UprightFont=*-regular,BoldFont=*-bold,ItalicFont=*-italic,BoldItalicFont=*-bolditalic]
\setmonofont{texgyrecursor}[Extension=.otf,UprightFont=*-regular,BoldFont=*-bold,ItalicFont=*-italic,BoldItalicFont=*-bolditalic,Scale=MatchLowercase]
\usepackage{amsmath,amssymb}
\usepackage{unicode-math}
\setmathfont{texgyrepagella-math.otf}
\usepackage{xcolor,microtype,enumitem,needspace,fancyhdr}
\usepackage{bookmark}
\definecolor{ink}{HTML}{17344B}
\definecolor{accent}{HTML}{197D80}
\definecolor{muted}{HTML}{596773}
\hypersetup{colorlinks=true,linkcolor=accent,urlcolor=accent,pdftitle={RA-16 - Super-exponential
decay of the zero-permanent
probability},pdfauthor={Open Problems in Numerical Linear Algebra}}
\urlstyle{same}
\setlength{\parindent}{0pt}
\setlength{\parskip}{5pt plus 1pt minus 1pt}
\linespread{1.04}
\setlength{\emergencystretch}{3em}
\widowpenalty=10000
\clubpenalty=10000
\displaywidowpenalty=10000
\allowdisplaybreaks[1]
\setlist{nosep,leftmargin=1.5em}
\providecommand{\tightlist}{\setlength{\itemsep}{0pt}\setlength{\parskip}{0pt}}
\providecommand{\pandocbounded}[1]{#1}
\pagestyle{fancy}
\fancyhf{}
\fancyhead[L]{\sffamily\footnotesize\color{muted}OPEN PROBLEMS IN NUMERICAL LINEAR ALGEBRA}
\fancyhead[R]{\sffamily\bfseries\footnotesize\color{accent}RA-16}
\fancyfoot[L]{\parbox[b]{0.9\textwidth}{\sffamily\fontsize{8}{10}\selectfont\color{muted}Editorial ratings; a bounded search cannot certify that a problem remains unsolved.\\Literature check: 2026-09-10}}
\fancyfoot[R]{\sffamily\footnotesize\color{muted}\thepage}
\renewcommand{\headrulewidth}{0.3pt}
\renewcommand{\headrule}{\hbox to\headwidth{\color{accent}\leaders\hrule height \headrulewidth\hfill}}
\makeatletter
\renewcommand{\section}{\Needspace{5\baselineskip}\@startsection{section}{1}{0pt}{12pt plus 2pt minus 2pt}{4pt}{\sffamily\bfseries\large\color{ink}}}
\renewcommand{\subsection}{\Needspace{4\baselineskip}\@startsection{subsection}{2}{0pt}{9pt plus 1pt minus 1pt}{3pt}{\sffamily\bfseries\normalsize\color{ink}}}
\makeatother
\setcounter{secnumdepth}{0}
\begin{document}
{\sffamily\small\color{muted}Randomized and low-rank approximation\par}
\vspace{3pt}
{\raggedright\hyphenpenalty=10000\exhyphenpenalty=10000\sffamily\bfseries\fontsize{21}{24}\selectfont\color{ink}Super-exponential
decay of the zero-permanent probability\par}
\vspace{7pt}
{\sffamily\small\textbf{Difficulty:} extreme\quad\textbf{Importance:} interesting
to the community\par}
{\sffamily\small\bfseries\color{accent}Status: Open\par}
\vspace{4pt}
{\color{accent}\hrule height 0.8pt}
\vspace{6pt}
\textbf{Rating rationale:} Extreme difficulty reflects the unresolved
passage from a fixed exponential rate to arbitrarily large rates,
despite the recent exponential breakthrough. The question concerns a
central cancellation phenomenon in combinatorial random-matrix theory,
supporting community importance.

\subsection{Problem statement}\label{problem-statement}

For each integer \(n\ge1\), let \(A_n=(a_{ij})\in\mathbb R^{n\times n}\)
have independent entries with \(\Pr(a_{ij}=1)=\Pr(a_{ij}=-1)=1/2\).
Define \[
\operatorname{per}A_n=\sum_{\sigma\in S_n}\prod_{i=1}^n a_{i,\sigma(i)}.
\] Is it true that, for every real \(c>0\), there exists \(N_c\) such
that \[
\Pr(\operatorname{per}A_n=0)\le e^{-cn}\qquad\text{for every }n\ge N_c?
\] Equivalently,
\(\lim_{n\to\infty}\Pr(\operatorname{per}A_n=0)^{1/n}=0\). This
explicitly quantifies the source's phrase ``super exponentially small.''

\subsection{Relevance and ratings}\label{relevance-and-ratings}

The question concerns cancellation in a multilinear function of a random
matrix. It complements random-matrix invertibility theory and bounds for
randomized permanent calculations, with connections to combinatorial
probability. The large gap beyond the current exponential estimate
justifies extreme difficulty.

\subsection{References}\label{references}

\begin{itemize}
\tightlist
\item
  V. H. Vu, \emph{Recent progress in combinatorial random matrix
  theory}, Probability Surveys 18 (2021), 179--200, Conjecture 6.12
  (\href{https://arxiv.org/pdf/2005.02797}{primary paper};
  \href{https://doi.org/10.1214/20-PS346}{journal}).
\item
  Z. Hunter, M. Kwan and L. Sauermann, \emph{Exponential
  anticoncentration of the permanent} (2025), Theorem 1.1 and §1.4
  (\href{https://arxiv.org/abs/2509.22577}{primary manuscript};
  \href{https://mkwn.github.io/EAP.pdf}{author PDF}).
\end{itemize}

\subsection{Status check}\label{status-check}

Hunter--Kwan--Sauermann prove the displayed bound for one absolute
positive constant, and explicitly identify Vu's stronger assertion as a
future direction. Their result resolves exponential growth of the
permanent range, which is not counted separately. Searches for Vu's
conjecture, super-exponential permanent anticoncentration, and
2025--2026 found no complete resolution. The July 2026
\href{https://arxiv.org/abs/2607.20329}{Ginibre-ensemble paper} treats
Gaussian entries and does not establish the discrete Bernoulli
assertion.

\subsection{Audit --- 2026-09-10}\label{audit-2026-09-10}

Independently rechecked \href{https://arxiv.org/pdf/2005.02797}{Vu,
Conjecture 6.12} and
\href{https://arxiv.org/html/2509.22577}{Hunter--Kwan--Sauermann,
Theorem 1.1 and §1.4}; the latter's arXiv record still lists only its
September 2025 version and explicitly leaves super-exponential decay
open. The fixed-rate theorem is a weaker bound, so the status remains
Open. Targeted later searches also checked the authors'
\href{https://arxiv.org/html/2603.15856}{March 2026 finite-field paper,
Theorems 1.3--1.4} and the \href{https://arxiv.org/abs/2607.20329}{July
2026 Gaussian paper}: their different field or entry distribution does
not prove this real Rademacher assertion. No complete resolution was
located in this bounded search.
\end{document}
